\documentclass[10pt]{article}
\usepackage[margin=1in]{geometry}
\usepackage{amsmath,amssymb}
\usepackage{booktabs}
\usepackage{graphicx}
\usepackage[hidelinks]{hyperref}
\usepackage[nohyphen]{underscore}

% Every number in this paper is a macro generated by make_numbers.py from numbers_map.json.
% Generated by make_numbers.py from numbers_map.json; do not edit.
\newcommand{\OpiPrime}{2305843009213693951}
\newcommand{\OpiN}{127}
\newcommand{\OpiKLow}{39}
\newcommand{\OpiKHigh}{77}
\newcommand{\OpiR}{2}
\newcommand{\OpiSeedA}{29}
\newcommand{\OpiSeedB}{71}
\newcommand{\OpiTrials}{8}
\newcommand{\OpiProposals}{8}
\newcommand{\OpiLowGlobalA}{108}
\newcommand{\OpiLowGlobalB}{110}
\newcommand{\OpiLowMixedInitA}{102}
\newcommand{\OpiLowMixedA}{109}
\newcommand{\OpiLowMixedInitB}{102}
\newcommand{\OpiLowMixedB}{109}
\newcommand{\OpiHighGlobalA}{127}
\newcommand{\OpiHighGlobalB}{127}
\newcommand{\OpiHighMixedInitA}{118}
\newcommand{\OpiHighMixedA}{124}
\newcommand{\OpiHighMixedInitB}{118}
\newcommand{\OpiHighMixedB}{122}
\newcommand{\NormScoreA}{127}
\newcommand{\NormScoreAStrict}{127}
\newcommand{\NormScoreB}{127}
\newcommand{\NormScoreBStrict}{127}
\newcommand{\NormDegA}{69}
\newcommand{\NormDegB}{64}
\newcommand{\NormTrialsA}{1}
\newcommand{\NormTrialsB}{15}
\newcommand{\NormTrialCap}{128}
\newcommand{\NormTargetDegA}{58}
\newcommand{\NormTargetDegB}{64}
\newcommand{\NormSuffDimA}{70}
\newcommand{\NormSuffDimB}{65}
\newcommand{\MorseSmallN}{15}
\newcommand{\MorseSmallK}{12}
\newcommand{\MorseSmallScore}{15}
\newcommand{\MorseSmallBound}{1138227279843745126/11529215046068469755}
\newcommand{\MorseCoverDeg}{76}
\newcommand{\MorseOrigBound}{0}
\newcommand{\MorseOrigBoundB}{0}
\newcommand{\MorseLiftScore}{127}
\newcommand{\MorseFiberBudget}{1672}
\newcommand{\MorseFailBits}{10}
\newcommand{\ExtOrigBound}{0}
\newcommand{\ExtOrigTwist}{0}
\newcommand{\ExtThreshold}{397893820763338319480932131383527086802587554693376}
\newcommand{\ExtLiftPrime}{6277101735386680763835789423207666416083908700390324961279}
\newcommand{\ExtLiftScore}{127}
\newcommand{\ExtLiftTwist}{0}
\newcommand{\ReadEllA}{19}
\newcommand{\ReadEllB}{22}
\newcommand{\ReadEllC}{24}
\newcommand{\ReadEllD}{32}
\newcommand{\ReadClassical}{109}
\newcommand{\ReadThreshA}{110}
\newcommand{\ReadThreshB}{115}
\newcommand{\ReadTailA}{0}
\newcommand{\ReadTailB}{0}
\newcommand{\ReadTailC}{0}
\newcommand{\DcpQ}{64}
\newcommand{\DcpHeaderA}{16}
\newcommand{\DcpHeaderB}{17}
\newcommand{\DcpSecret}{1}
\newcommand{\DcpNoteX}{0.7745}
\newcommand{\DcpNoteLocal}{0.7410}
\newcommand{\DcpNoteColl}{0.8801}
\newcommand{\DcpTwoPackets}{2}
\newcommand{\DcpThreePackets}{3}
\newcommand{\CsSeed}{0}
\newcommand{\CsRestarts}{12}
\newcommand{\CsModGamma}{1.0}
\newcommand{\CsTwoGammaA}{0.85}
\newcommand{\CsTwoGammaB}{0.8}
\newcommand{\CsThreeGammaA}{0.85}
\newcommand{\CsThreeGammaB}{0.8}
\newcommand{\CsThreeGammaC}{0.75}
\newcommand{\CsThreeKnownCircuits}{487,821}
\newcommand{\CsModKnownX}{0.7745}
\newcommand{\CsModKnownLocHel}{0.7410}
\newcommand{\CsModKnownProd}{0.8483}
\newcommand{\CsModKnownAdapt}{0.8801}
\newcommand{\CsModKnownPPT}{0.8801}
\newcommand{\CsModKnownColl}{0.8801}
\newcommand{\CsModKnownGap}{0.0000}
\newcommand{\CsTwoKnownX}{0.4763}
\newcommand{\CsTwoKnownLocHel}{0.5657}
\newcommand{\CsTwoKnownProd}{0.5895}
\newcommand{\CsTwoKnownAdapt}{0.6147}
\newcommand{\CsTwoKnownPPT}{0.6147}
\newcommand{\CsTwoKnownColl}{0.6147}
\newcommand{\CsTwoKnownGap}{0.0000}
\newcommand{\CsThreeKnownX}{0.5746}
\newcommand{\CsThreeKnownLocHel}{0.6495}
\newcommand{\CsThreeKnownProd}{0.6698}
\newcommand{\CsThreeKnownAdapt}{0.6868}
\newcommand{\CsThreeKnownPPT}{0.6895}
\newcommand{\CsThreeKnownColl}{0.6993}
\newcommand{\CsThreeKnownGap}{0.0098}
\newcommand{\CsThreeKnownWindow}{0.0027}
\newcommand{\CsModUnkX}{0.7745}
\newcommand{\CsModUnkProd}{0.7745}
\newcommand{\CsModUnkAdapt}{0.7745}
\newcommand{\CsModUnkPPT}{0.8305}
\newcommand{\CsModUnkColl}{0.8305}
\newcommand{\CsModUnkGap}{0.0000}
\newcommand{\CsModUnkWindow}{0.0560}
\newcommand{\CsTwoUnkX}{0.4763}
\newcommand{\CsTwoUnkProd}{0.4763}
\newcommand{\CsTwoUnkAdapt}{0.4967}
\newcommand{\CsTwoUnkPPT}{0.5122}
\newcommand{\CsTwoUnkColl}{0.5143}
\newcommand{\CsTwoUnkGap}{0.0022}
\newcommand{\CsTwoUnkWindow}{0.0155}
\newcommand{\CsThreeUnkX}{0.5746}
\newcommand{\CsThreeUnkProd}{0.5783}
\newcommand{\CsThreeUnkAdapt}{0.6263}
\newcommand{\CsThreeUnkPPT}{0.6263}
\newcommand{\CsThreeUnkColl}{0.6369}
\newcommand{\CsThreeUnkGap}{0.0106}
\newcommand{\CsThreeUnkWindow}{4.0\times 10^{-6}}
\newcommand{\PolarN}{32768}
\newcommand{\PolarF}{16384}
\newcommand{\PolarRate}{1/2}
\newcommand{\PolarOverlap}{13/20}
\newcommand{\PolarSourceP}{7/40}
\newcommand{\PolarTV}{424553765284920852120940002337825/41538374868278621028243970633760768}
\newcommand{\PolarTarget}{1/10}
\newcommand{\PolarMaxF}{17027}
\newcommand{\PolarMaxRate}{17027/32768}
\newcommand{\PolarGrid}{256}
\newcommand{\PolarPower}{3}
\newcommand{\PolarMassBits}{96}
\newcommand{\PolarCombos}{65,534}
\newcommand{\PolarVisits}{1,008,335,816}
\newcommand{\CoarseGrid}{64}
\newcommand{\CoarsePower}{4}
\newcommand{\CoarseNA}{64}
\newcommand{\CoarseFA}{16}
\newcommand{\CoarseRateA}{1/4}
\newcommand{\CoarseHalfA}{1}
\newcommand{\CoarseNB}{256}
\newcommand{\CoarseFB}{78}
\newcommand{\CoarseRateB}{39/128}
\newcommand{\CoarseHalfB}{1}
\newcommand{\CoarseNC}{1024}
\newcommand{\CoarseFC}{361}
\newcommand{\CoarseRateC}{361/1024}
\newcommand{\CoarseHalfC}{1}
\newcommand{\CoarseND}{4096}
\newcommand{\CoarseFD}{1614}
\newcommand{\CoarseRateD}{807/2048}
\newcommand{\CoarseHalfD}{1}
\newcommand{\EighthN}{4096}
\newcommand{\EighthRate}{3/8}
\newcommand{\CalN}{2}
\newcommand{\CalNoise}{1/10}
\newcommand{\CalAvg}{2/5}
\newcommand{\CalWorst}{18/25}
\newcommand{\RedTotal}{976}
\newcommand{\RedPublished}{309}
\newcommand{\RedAdjacent}{452}
\newcommand{\RedNotFound}{2}
\newcommand{\RedUnclear}{213}
\newcommand{\RedShown}{366}
\newcommand{\RedPropagated}{610}
\newcommand{\RedGroundN}{10}
\newcommand{\RedGroundPub}{7}
\newcommand{\RedSpotAgree}{27}
\newcommand{\RedSpotN}{30}
\newcommand{\RedSpotShownAgree}{11}
\newcommand{\RedSpotShownN}{12}
\newcommand{\RedSpotPropAgree}{16}
\newcommand{\RedSpotPropN}{18}
\newcommand{\RedCodePub}{168}
\newcommand{\RedCodeAdj}{141}
\newcommand{\RedCodeUnc}{148}
\newcommand{\RedCosetPub}{121}
\newcommand{\RedCosetAdj}{115}
\newcommand{\RedCosetNf}{1}
\newcommand{\RedCosetUnc}{17}
\newcommand{\RedDcpPub}{11}
\newcommand{\RedDcpAdj}{111}
\newcommand{\RedDcpNf}{1}
\newcommand{\RedDcpUnc}{31}
\newcommand{\RedToyAdj}{53}
\newcommand{\RedQueryPub}{1}
\newcommand{\RedQueryAdj}{19}
\newcommand{\RedQueryUnc}{2}
\newcommand{\RedOtherPub}{2}
\newcommand{\RedOtherAdj}{9}
\newcommand{\RedOtherUnc}{6}
\newcommand{\RedAppliedPub}{5}
\newcommand{\RedAppliedUnc}{6}
\newcommand{\RedLatticePub}{1}
\newcommand{\RedLatticeAdj}{4}
\newcommand{\RedLatticeUnc}{2}
\newcommand{\RedDqiUnc}{1}
\newcommand{\FmTotal}{976}
\newcommand{\FmClassical}{311}
\newcommand{\FmProofGap}{149}
\newcommand{\FmNoSignal}{116}
\newcommand{\FmReadout}{95}
\newcommand{\FmKappa}{0.86}
\newcommand{\FmSecondN}{214}
\newcommand{\OpiPrimeBits}{61}
\newcommand{\ExtThresholdBits}{169}
\newcommand{\TwistThresholdBits}{756}
\newcommand{\ExtLiftPrimeBits}{192}
\newcommand{\MorseLiftPrimeBits}{1279}
\newcommand{\MorseSmallBoundDec}{0.0987}
\newcommand{\ReadMeanA}{0.8169}
\newcommand{\ReadMeanB}{0.8406}
\newcommand{\ReadMeanC}{0.8550}
\newcommand{\ReadMeanD}{0.9028}
\newcommand{\ReadTailDA}{0.3145}
\newcommand{\ReadTailDB}{0.0509}
\newcommand{\CubeN}{63}
\newcommand{\CubeK}{46}
\newcommand{\CubeR}{3}
\newcommand{\CubeMixedA}{55}
\newcommand{\CubeMixedB}{56}
\newcommand{\QuartN}{32}
\newcommand{\QuartK}{24}
\newcommand{\QuartR}{4}
\newcommand{\QuartPrime}{998244353}
\newcommand{\QuartMixedA}{28}
\newcommand{\QuartMixedB}{27}
\newcommand{\CsModUnkPrimal}{0.83052061}
\newcommand{\CsModUnkCollLong}{0.83052070}
\newcommand{\CsModThetaA}{\pi/2}
\newcommand{\CsModThetaB}{17\pi/32}
\newcommand{\CsTwoThetaA}{\pi/3}
\newcommand{\CsTwoThetaB}{\pi/2}
\newcommand{\CsThreeThetaA}{\pi/4}
\newcommand{\CsThreeThetaB}{\pi/3}
\newcommand{\CsThreeThetaC}{2\pi/3}
\newcommand{\PolarTVDec}{0.010221}


\newcommand{\TV}{\mathrm{TV}}
\newcommand{\Fp}{\mathbb{F}_p}

\title{Classical competitors and local measurements for three proposed quantum readouts and samplers}
\author{Q-Search project (draft, not for circulation)}
\date{}

\begin{document}
\maketitle

\begin{abstract}
The Q-Search project proposes quantum mechanisms, attacks each one classically, and keeps the negative results. This draft reports three cases in which a proposed quantum readout or sampler meets a classical competitor, or a local measurement that matches it.
(A)~For explicit prime-field quadratic-character instances of optimal polynomial intersection (OPI) with $N=\OpiN$ points over $p=\OpiPrime$, classical algebraic attacks reach perfect $\NormScoreA/\OpiN$ witnesses on every saved input with $K=\OpiKHigh$. In the original field the available fixed-input lower bounds for a Morse-theoretic solver are zero, which is not evidence of hardness; a positive fixed-input bound was obtained only for one input re-posed over larger declared primes, which is a different problem. At the low rate $K=\OpiKLow$, on the mixed input that a finite quantum readout certificate addresses, the best classical witness is $\OpiLowMixedA/\OpiN$; the certificate exceeds it only under an exact source and an exact unique decoder that has not been built, so this case is open.
(B)~For binary readout of multi-packet dihedral coset problem (DCP) phase packets, on the recorded known-secret pure instance a two-step adaptive local measurement with one bit of feed-forward reaches the collective Helstrom score $\CsModKnownColl$. This is what the theorem of Virmani, Sacchi, Plenio and Markham predicts. The collective measurement beats the two product rules the bridge note evaluated and the best product receiver found by search, but not the best adaptive local strategy. On the three-packet mixed instances a certified PPT bound leaves gaps of at least $\CsThreeKnownGap$ and $\CsThreeUnkGap$ between every local protocol and the collective optimum; on the two-packet mixed instances the adaptive search reaches the collective optimum with a known secret, and with unknown sign the certified gap is below $\CsTwoUnkGapUp$. The unknown-sign case of the recorded instance is open: the best local protocol lies between $\CsModUnkAdapt$ and $\CsModUnkPPT$.
(C)~For signed polar-code Gibbs sampling with uniformly random signs, a classical successive-cancellation source sampler with a certified public design has average total-variation error at most $\PolarTVDec$ (an exact rational bound, rounded up) at length $N=\PolarN$ and rate $\PolarRate$. The certificate is conditional on review-pending local derivations, among them the correspondence between posterior magnitudes and pure-state channels and the degradation argument that makes each computed term an upper bound. It covers the uniform-sign average only.
Separately, a rediscovery audit of the project's $\RedTotal$ retained negative results, which does not cover Results A--C, labels $\RedPublished$ of them as restating a published barrier; $\RedReadPublished$ of these are among the $\RedShown$ records a judge read, and the rest inherit their cluster's label. No result here is a new quantum algorithm, a speedup claim or a hardness claim, and most derivations are marked review pending in their source notes.
\end{abstract}

\section{Introduction}
\label{sec:intro}

The Q-Search project runs a fixed loop. It proposes a quantum mechanism for an algorithmic task, attacks it with the strongest classical algorithm or restricted measurement it can find, and records the outcome. Failed proposals are kept, not discarded: the negative-results registry holds $\RedTotal$ records. A failure-mode meta-analysis of that registry assigns each record one primary failure mode. ``Classical competitor'' is the most common, with $\FmClassical$ of $\FmTotal$ records, ahead of proof gaps ($\FmProofGap$), absent signal ($\FmNoSignal$) and readout cost ($\FmReadout$). A second run of the same model, in a fresh context, relabelled $\FmSecondN$ sampled groups ($\FmSecondRecords$ records) and agreed with the first run per group with Cohen's $\kappa=\FmKappa$. That measures how repeatable one model's labels are, not agreement with an independent judge.

This draft collects three recent instances of the loop. In each, a quantum readout or sampler was proposed, and the comparison that decides its value turned out to be classical, or local:
\begin{itemize}
\item \textbf{A. Prime-character OPI} (Section~\ref{sec:opi}). Classical power completion and a quadratic-extension norm search find perfect witnesses on every saved $N=\OpiN$, $K=\OpiKHigh$ input with explicit quadratic-character lists; one smaller saved quadratic input is not solved perfectly (Section~\ref{sec:limits}). A Morse-theoretic argument gives a fixed-input guarantee only in an explicit large-field regime, which the saved $N=\OpiN$ inputs over the $\OpiPrimeBits$-bit prime do not meet. The low-rate instance remains open.
\item \textbf{B. Multi-packet DCP readout} (Section~\ref{sec:dcp}). A circuit search over local adaptive receivers, together with a certified PPT bound, decides on some instances whether a reported collective-measurement advantage survives against local strategies, and brackets the answer on the others.
\item \textbf{C. Polar Gibbs source} (Section~\ref{sec:polar}). A classical successive-cancellation (SC) sampler with a public frozen set, certified under review-pending local derivations, matches the uniform-sign Gibbs task that a natural polar quantum receiver would target.
\end{itemize}
Section~\ref{sec:lit} summarizes what a rediscovery audit of the registry says the literature already had; it does not cover Results A--C. Section~\ref{sec:limits} lists limitations, and Section~\ref{sec:repro} explains how the numbers are regenerated.

\paragraph{Conventions.} Every number in the text and tables is a macro read from a committed JSON report (Section~\ref{sec:repro}). Where a report stores a quantity only as an exact rational, we print that rational, a directed decimal or a qualitative description; we never round by hand. The script rounds directed decimals the way that keeps each statement true: down for lower bounds and for scores that a measurement or witness achieves, up for upper bounds and for the widths of open windows. The only decimal rounded to nearest is Cohen's $\kappa$. Quotations in the ``Scope'' paragraphs are verbatim from the source notes or reports. The source notes mark most local derivations as review pending, and none of the results below has had independent review.

%--------------------------------------------------------------------
\section{Result A: a classical competitor for prime-character OPI}
\label{sec:opi}

\paragraph{Setting.} Optimal polynomial intersection (OPI) asks for a polynomial of degree less than $K$ whose values at $N$ points land in as many prescribed sets as possible. It is a central target problem of decoded quantum interferometry (DQI)~\cite{JordanDQI2024}. Sun and Wootters give a quantum algorithm for OPI over prime fields whose asymptotic guarantee approaches a semicircle-law target and becomes perfect search at fixed rate above a threshold~\cite{SunWootters2026}. Their guarantee is asymptotic, and their oracle separation concerns random membership tables.

The project instantiates OPI with explicit sets. Fix a prime $p$, the points $1,\ldots,N$ in $\Fp$, a power index $r \mid p-1$, public shifts $t_i$, orientations $\sigma_i$ and a verified coset generator $g$. Position $i$ accepts $Q$ when
\[
Q(i) - t_i \in g^{\sigma_i} \cdot (\Fp^{\ast})^{r} \quad (\text{optionally also } Q(i)-t_i=0),
\]
and the task is to maximize the number of accepted positions over $\deg Q < K$. Zero is included in one mode and excluded in the ``strict'' mode. ``Global'' inputs use a common coset; ``mixed'' inputs use seeded independent orientations. Membership is decided by modular exponentiation: no acceptance table and no oracle is supplied. The saved large instances have $p=\OpiPrime$, $N=\OpiN$, $r=\OpiR$, public seeds $\OpiSeedA$ and $\OpiSeedB$, and dimension $K=\OpiKHigh$ (called high rate in the notes) or $K=\OpiKLow$ (low rate).

\paragraph{Claim tested.} The Morse note states the opening plainly: ``The prospective quantum advantage in this family was polynomial perfect intersection at high rate.'' The question was whether explicit power-character predicates keep perfect high-rate intersection out of classical reach, without hiding a random membership oracle.

\paragraph{Classical competitors.} Four attacks were run, each returning native polynomial witnesses that are rechecked by independent Horner evaluation.
\begin{enumerate}
\item \emph{Power completion.} For $M$ positions in a common coset, a formal $r$-th root at infinity yields $Q$ with the required residues. Every common-coset input is solved perfectly whenever $K \ge (r-1)\lceil N/r \rceil$; the strict mode needs one extra degree.
\item \emph{Quadratic jet with adaptive root-value exchange} for mixed cosets. Free root coefficients are perturbed through Lagrange basis polynomials, and only strict improvements of the whole native word are kept. Each source pays $\OpiTrials$ trials per mechanism and $\OpiProposals$ proposals per coordinate. Runs stop with \texttt{NO\_SAMPLED\_IMPROVEMENT}, which is not a certified local optimum.
\item \emph{Quadratic-extension norm search.} Write the residual target as a norm $F = R^2 - \nu Z^2$ from $\mathbb{F}_{p^2}$, which exists exactly when every odd-degree irreducible factor has even multiplicity; then $Q$ is read off from $R$. Strict roots are obtained by a norm-one Cayley rotation. The splitting and norm facts are standard properties of constant field extensions~\cite{Rosen2002}; for $\nu=-1$, a nonsquare here because $p \equiv 3 \pmod 4$, see also~\cite{KuoLiuTotani2026}. The method needs the degree budget $K > \max(|A|, D)$, where $A$ is one orientation class and $D$ the padded size of the other. Under independent uniform shifts the first trial succeeds with probability at least $I_p(D)/p^D$, where $I_p(D)$ counts monic irreducibles of degree $D$; this is a population bound, not a fixed-input one.
\item \emph{Morse progression with a fixed-input fibre bound.} A public arithmetic progression of norm targets $f - v g$ is built; the code checks exact tame Morse conditions, which give monodromy $S_n$ under the note's argument. A Weil bound on a twisted Galois closure~\cite{BankBarySorokerRosenzweig2013}, or a sharper exterior-power trace bound~\cite{FouvryKowalskiMichel2014}, then lower-bounds the fraction of irreducible fibres for that fixed input. Each bound is guaranteed to give a per-fibre rate of at least $1/(2n)$, where $n$ is the degree of the cover, once $p$ is at least an explicit threshold.
\end{enumerate}

\begin{table}[t]
\centering
\small
\caption{Result A, native scores on the saved instances ($p=\OpiPrime$, $N=\OpiN$, $r=\OpiR$, zero included unless stated). An arrow shows the adaptive repair from its starting word. Sources: \texttt{prime\_character\_opi.json} and \texttt{prime\_character\_norm\_attack.json} in \texttt{research/classical\_baselines/}.}
\label{tab:opi}
\begin{tabular}{llll}
\toprule
$K$ & Input & Attack & Score (of $\OpiN$) \\
\midrule
\OpiKHigh & global, seeds \OpiSeedA{} / \OpiSeedB & power completion & \OpiHighGlobalA{} / \OpiHighGlobalB \\
\OpiKHigh & mixed, seed \OpiSeedA & jet, adaptive repair & \OpiHighMixedInitA{} $\to$ \OpiHighMixedA \\
\OpiKHigh & mixed, seed \OpiSeedB & jet, adaptive repair & \OpiHighMixedInitB{} $\to$ \OpiHighMixedB \\
\OpiKHigh & mixed, seed \OpiSeedA & norm search (zero incl.\ / strict) & \NormScoreA{} / \NormScoreAStrict \\
\OpiKHigh & mixed, seed \OpiSeedB & norm search (zero incl.\ / strict) & \NormScoreB{} / \NormScoreBStrict \\
\midrule
\OpiKLow & global, seeds \OpiSeedA{} / \OpiSeedB & power completion & \OpiLowGlobalA{} / \OpiLowGlobalB \\
\OpiKLow & mixed, seed \OpiSeedA & jet, adaptive repair & \OpiLowMixedInitA{} $\to$ \OpiLowMixedA \\
\OpiKLow & mixed, seed \OpiSeedB & jet, adaptive repair & \OpiLowMixedInitB{} $\to$ \OpiLowMixedB \\
\OpiKLow & mixed, both seeds & norm search & not applicable \\
\bottomrule
\end{tabular}
\end{table}

\paragraph{Results.} Table~\ref{tab:opi} gives the scores. At $K=\OpiKHigh$ every saved input has a perfect classical witness. The global inputs fall to power completion. The mixed inputs stop at $\OpiHighMixedA$ and $\OpiHighMixedB$ under adaptive repair, and the norm search then solves both perfectly, in the zero-including and strict modes, with the shifts and orientations unchanged. The witnesses have degree $\NormDegA$ and $\NormDegB$, inside the dimension budget, and were found after $\NormTrialsA$ and $\NormTrialsB$ charged base-field factorizations (trial cap $\NormTrialCap$). The exact population lower bounds use norm-target degrees $\NormTargetDegA$ and $\NormTargetDegB$; the norm note describes them as approximately $1/D$ and warns that they are not probabilities for these particular seeded shifts. At $K=\OpiKLow$ the norm search does not apply: the degree budget asks for dimension $\NormSuffDimA$ and $\NormSuffDimB$.

\begin{table}[t]
\centering
\small
\caption{Result A, fixed-input lower bounds on the fraction of irreducible fibres for the Morse norm solver (quadratic, mixed orientations). ``Original'' is $p=\OpiPrime$. $^\dagger$Zero included; the twist bound on seed $\OpiSeedB$ is also $\MorseOrigBoundB$. A zero lower bound is not a zero probability. Sources: \texttt{prime\_character\_morse\_attack.json} and \texttt{prime\_character\_exterior\_bound.json}.}
\label{tab:morse}
\begin{tabular}{llllll}
\toprule
$(N, K)$ & Input & Field & Twist bound & Exterior bound & Witness \\
\midrule
(\MorseSmallN, \MorseSmallK) & controls, both zero modes & original & positive & not evaluated & \MorseSmallScore{} of \MorseSmallN \\
(\OpiN, \OpiKHigh) & mixed, seed \OpiSeedA{}$^\dagger$ & original & \MorseOrigBound{} & \ExtOrigBound{} & none from this route \\
(\OpiN, \OpiKHigh) & strict, seed \OpiSeedA & $\MorseLiftPrimeBits$-bit prime & positive & not evaluated & \MorseLiftScore{} of \OpiN \\
(\OpiN, \OpiKHigh) & strict, seed \OpiSeedA & Eq.~(\ref{eq:primes}) & \ExtLiftTwist{} & positive & \ExtLiftScore{} of \OpiN \\
\bottomrule
\end{tabular}
\end{table}

The Morse and exterior bounds answer a different question: is there a guarantee for a fixed input, not just a seeded success or a population average? Table~\ref{tab:morse} shows the answer depends on the field. In the original field the small controls ($N=\MorseSmallN$, $K=\MorseSmallK$) have a positive certified bound; for the first control it is
\[
\MorseSmallBound \;\ge\; \MorseSmallBoundDec
\]
per fibre. For the large $K=\OpiKHigh$ inputs the geometry is verified, but both the twist bound ($\MorseOrigBound$) and the exterior bound ($\ExtOrigBound$) are zero: in each, the error term exceeds $p$. A per-fibre rate of at least $1/(2n)$, with $n=\MorseCoverDeg$ the cover degree, is guaranteed once $p$ is at least the bound's threshold: a $\TwistThresholdBits$-bit integer for the twist bound, and for the exterior bound the $\ExtThresholdBits$-bit integer $p_{\text{threshold}}$ of Eq.~(\ref{eq:threshold}). A prime with as many bits as a threshold can still lie below it; every prime of at least $\TwistThresholdSafeBits$ bits meets the twist threshold, and every prime of at least $\ExtThresholdSafeBits$ bits meets the exterior one, against $\OpiPrimeBits$ bits for $p$. A zero lower bound says only that these conservative bounds do not apply; it is not evidence of hardness. In a declared larger prime, with the same shifts and orientations of the strict seed-$\OpiSeedA$ input, the exterior bound is positive while the twist bound is $\ExtLiftTwist$, and a perfect strict witness is found:
\begin{equation}
\label{eq:threshold}
p_{\text{threshold}} = \ExtThreshold,
\end{equation}
\begin{equation}
\label{eq:primes}
p_{\text{declared}} = \ExtLiftPrime \quad (\ExtLiftPrimeBits \text{ bits}).
\end{equation}
Here $p \ge p_{\text{threshold}}$ is the exterior method's sufficient condition, which the original prime does not meet and the declared prime does.
The large run had a budget of $\MorseFiberBudget$ fibres, which bounds failure by $2^{-\MorseFailBits}$ under independent uniform fibre coins; the seeded run is a replay, not a confidence certificate. A lift to another field changes the problem, so these witnesses do not solve the original input.

\begin{table}[t]
\centering
\small
\caption{Result A, low rate ($K=\OpiKLow$, $N=\OpiN$): finite character-DQI readout certificates, conditional on an exact source and an exact verified unique decoder, against the saved same-input classical witness $\ReadClassical/\OpiN$. The certified mean is the expected fraction of the $\OpiN$ positions satisfied, printed as an upper bound (rounded up) where it falls below $\ReadClassical/\OpiN$ and as a lower bound (rounded down) where it exceeds it; the report stores exact rationals. A zero lower bound on the tail is not a zero probability. Source: \texttt{prime\_character\_finite\_readout.json}.}
\label{tab:readout}
\begin{tabular}{llll}
\toprule
Radius $\ell$ & Mean certificate & Certified mean & One-attempt lower bound on reaching $\ReadThreshA$ ($\ReadThreshB$) \\
\midrule
\ReadEllA & exact RS moments & $\le \ReadMeanA$ (below) & \ReadTailA \\
\ReadEllB & extended radial Gram & $\le \ReadMeanB$ (below) & \ReadTailB \\
\ReadEllC & extended radial Gram & $\le \ReadMeanC$ (below) & \ReadTailC \\
\ReadEllD & extended radial Gram & $\ge \ReadMeanD$ (above) & $\ReadTailDA$ ($\ReadTailDB$) \\
\bottomrule
\end{tabular}
\end{table}

At $K=\OpiKLow$ the comparison runs the other way, and the case is open. The best classical witness is $\OpiLowMixedA/\OpiN$ on both mixed seeds, from adaptive repair; it is a witnessed score, not an optimum and not a limit on what classical methods can reach. On the quantum side, a known DQI route (sparse Fourier errors, a public Reed--Solomon syndrome, coherent unique decoding, Fourier readout) has a finite certificate for the mixed seed-$\OpiSeedA$ input. Among the radii evaluated, its certified mean exceeds the classical witness only at $\ell=\ReadEllD$ (Table~\ref{tab:readout}), and only conditionally: the certificate assumes an exact source and an exact unique decoder beyond half the minimum distance. That condition carries the whole comparison. Mazumder and Niroula prove that for the satisfied-count objective under worst-case Hamming decoding ($2\ell+1<d$), DQI's satisfied fraction exceeds Prange completion, $\rho+(1-\rho)\alpha$, by at most $\sqrt{u(1-u)}-u$ with $u=\alpha/2$, and never by more than $(\sqrt2-1)/2$~\cite{MazumderNiroula2026}. With code rate $\alpha=\OpiKLow/\OpiN$ and accepted density $\rho\approx 1/2$ this caps the guarantee near $0.861$ of the positions, which the classical witness already nearly reaches. The bound is asymptotic, so at $N=\OpiN$ it is a guide, not a certificate. The fixed-gap decoder the note audits, by Dao, Kominers and Thaler~\cite{DaoKominersThaler2026}, has a gap-dependent exponent, and its generic finite certificate has a length threshold that the note finds cannot be applied directly at this length. Jeronimo's capacity list decoder~\cite{Jeronimo2026} depends polynomially on the field size, which the OPI note judges insufficient when $\log p$ is the charged input size. No decoder circuit has been built. The flat Fourier magnitudes that the source relies on use the quadratic Gauss sum~\cite{MurtyPathak2017} and need $p \equiv 3 \pmod 4$, which holds here. Neither side has a matched runtime comparison.

\paragraph{Scope.} The notes limit these results as follows. OPI note: ``Local derivations REVIEW PENDING. No novelty, hardness or speedup claim.'' and ``Do not call this an all-rate OPI dequantization.'' and ``This weakens the proposed mixed-coset escape, not the published quantum algorithm: a finite near-perfect classical word does not disprove asymptotic perfect quantum search.'' Norm note: ``Standard polynomial-factorization/norm machinery; local application REVIEW PENDING. No novelty, all-input dequantization or new quantum algorithm is claimed.'' and ``That is a population theorem, NOT a success bound for each fixed input or PRNG seed.'' and ``Saved K39 siblings are NOT APPLICABLE, not declared hard or unsatisfiable.'' Morse note: ``This is not a dequantization of arbitrary oracle lists, low-rate OPI, all fields, or every orientation pattern.'' and, on the original large inputs, ``The implementation refuses to promote their geometry to a probability claim. This is not evidence of classical hardness.'' Exterior-bound note: ``A lift to a different field does not solve or give a probability guarantee for that original input. No field-descent claim is made.'' Capacity-source note: ``But negligible ambiguity is NOT an efficient decoder.'' Finite-readout note: ``CONDITIONAL on an exact source and exact verified unique decoder.'' and ``It would be misleading to report this finite score certificate as an implemented quantum advantage or to infer classical hardness from it.''
\begin{flushleft}
The reports set these flags to \texttt{false}: \texttt{accepted\_speedup\_candidate}, \texttt{new\_quantum\_algorithm\_claimed} and \texttt{mixed\_cosets\_proved\_classically\_hard} (OPI report); \texttt{all\_input\_or\_asymptotic\_dequantization\_proved} (norm report); \texttt{all\_input\_or\_all\_field\_dequantization\_claimed} (Morse report); \texttt{all\_field\_or\_low\_rate\_dequantization\_claimed} (exterior report); and \texttt{new\_quantum\_algorithm\_or\_speedup\_claimed} and \texttt{matched\_classical\_quantum\_runtime\_comparison\_supplied} (finite-readout report).
\end{flushleft}

%--------------------------------------------------------------------
\section{Result B: local readout of multi-packet DCP phase packets}
\label{sec:dcp}

\paragraph{Setting.} In the dihedral coset problem (DCP)~\cite{Regev2004,Kuperberg2005}, each sample yields a qubit carrying a phase that depends on a public header and the hidden secret. Following the bridge note and the circuit search, a packet is modelled as
\[
\rho(\theta,\gamma) = \frac{1}{2}\left(I + \gamma\cos\theta\, X + \gamma\sin\theta\, Y\right),
\]
where $\theta = 2\pi h s/q$ for header $h$, secret $s$ and modulus $q$, and $\gamma=1$ for a pure packet. The task is binary: secret $0$, under which every packet has $\theta=0$ with the same $\gamma$, against secret $s^\ast$. Either $s^\ast$ is known, or its sign is unknown and the alternative hypothesis is the equal mixture of $s^\ast$ and $-s^\ast$, with all packet angles flipping sign together. A receiver is scored by the total-variation distance between its outcome laws under the two hypotheses; with equal priors and the optimal decision rule its success probability is $(1+\TV)/2$. The recorded instance has modulus $q=\DcpQ$, headers $\DcpHeaderA$ and $\DcpHeaderB$, secret $s^\ast=\DcpSecret$ and pure packets ($\gamma=\CsModGamma$), so its angles are $\CsModThetaA$ and $\CsModThetaB$. Two further instances, the bridge note's illustrative distinct-header controls, are given directly by their angles and use mixed packets: $\DcpTwoPackets$ packets with $\theta=(\CsTwoThetaA,\CsTwoThetaB)$ and $\gamma=(\CsTwoGammaA,\CsTwoGammaB)$, and $\DcpThreePackets$ packets with $\theta=(\CsThreeThetaA,\CsThreeThetaB,\CsThreeThetaC)$ and $\gamma=(\CsThreeGammaA,\CsThreeGammaB,\CsThreeGammaC)$.

\paragraph{Claim tested.} On the recorded instance the bridge note evaluates three readouts: an $X$-basis product readout ($\DcpNoteX$), a product of local Helstrom readouts ($\DcpNoteLocal$), and the collective Helstrom optimum over all joint measurements ($\DcpNoteColl$)~\cite{Helstrom1976}. The note withdraws an earlier universal hierarchy among them and states that it does not compare against the optimal adaptive or separable receiver. The gap between the collective score and the product scores reads as an advantage for an entangled measurement. The question is whether that advantage holds against the best local strategy, or only against the two product rules.

\paragraph{Local measurements and bounds.} A circuit search (\url{tools/circuit_search/dcp_readout.py}) supplies the missing comparison. Every receiver uses one single-qubit rotation and a $Z$ measurement per wire, and no entangling gate; adaptive receivers add classical feed-forward. \emph{Product} receivers measure each wire along its own optimized direction. \emph{Adaptive} receivers measure wires in sequence, each direction depending on earlier outcomes, over all wire orders. Both are optimized by Nelder--Mead and Powell from $\CsRestarts$ seeded restarts (seed $\CsSeed$); the three-packet known-secret search alone made \CsThreeKnownCircuits{} objective evaluations, each scoring one circuit. Two quantities are bounds rather than search results: the collective (Helstrom) optimum, computed in closed form as half the trace norm of the difference of the two hypothesis states, and a \emph{PPT bound}, computed by a semidefinite program as the best measurement whose effects stay positive under partial transpose across every cut. Every local protocol, including two-way, multi-round and POVM-based ones, is such a measurement, so the PPT value bounds all of them. The bound is certified by repairing the solver's dual into a feasible dual solution and adding the residual eigenvalue, so it rests on floating-point eigenvalue accuracy rather than on solver tolerance. The \emph{certified gap} is the collective score minus the PPT bound.

\begin{table}[t]
\centering
\caption{Result B, total-variation scores; $M$ is the number of packets and ``local Helstrom'' the product of local Helstrom readouts. ``Best product'' and ``best adaptive'' are search results, hence lower bounds on their optima. ``Collective'' is the Helstrom optimum. ``PPT bound'' is the smaller of the certified dual bound and the collective value; it bounds every local protocol. ``Certified gap'' is the collective value minus the PPT bound, a lower bound on the collective advantage over every local protocol. A certified gap of $\CsModKnownGap$ means that no gap was certified, not that none exists. Values are printed to $10^{-4}$ in the direction that keeps them valid: the PPT bound is rounded up, and every other entry (achieved scores, the collective value and the certified gap) is rounded down. Equal values can therefore print one unit apart: in the first, second and fourth rows the PPT bound equals the collective value, and in the first two rows the adaptive search reaches it. Gaps are computed from unrounded values. The local Helstrom product is defined only for a known secret. Source: \texttt{research/audits/circuit\_search\_dcp\_readout.json}.}
\label{tab:dcp}
\resizebox{\textwidth}{!}{%
\begin{tabular}{llcccccccc}
\toprule
Secret & Instance & $M$ & $X$ basis & Local Helstrom & Best product & Best adaptive & PPT bound & Collective & Certified gap \\
\midrule
known & mod-$\DcpQ$, pure & \DcpTwoPackets & \CsModKnownX & \CsModKnownLocHel & \CsModKnownProd & \CsModKnownAdapt & \CsModKnownPPT & \CsModKnownColl & \CsModKnownGap \\
known & mixed & \DcpTwoPackets & \CsTwoKnownX & \CsTwoKnownLocHel & \CsTwoKnownProd & \CsTwoKnownAdapt & \CsTwoKnownPPT & \CsTwoKnownColl & \CsTwoKnownGap \\
known & mixed & \DcpThreePackets & \CsThreeKnownX & \CsThreeKnownLocHel & \CsThreeKnownProd & \CsThreeKnownAdapt & \CsThreeKnownPPT & \CsThreeKnownColl & \CsThreeKnownGap \\
\midrule
unknown sign & mod-$\DcpQ$, pure & \DcpTwoPackets & \CsModUnkX & -- & \CsModUnkProd & \CsModUnkAdapt & \CsModUnkPPT & \CsModUnkColl & \CsModUnkGap \\
unknown sign & mixed & \DcpTwoPackets & \CsTwoUnkX & -- & \CsTwoUnkProd & \CsTwoUnkAdapt & \CsTwoUnkPPT & \CsTwoUnkColl & \CsTwoUnkGap \\
unknown sign & mixed & \DcpThreePackets & \CsThreeUnkX & -- & \CsThreeUnkProd & \CsThreeUnkAdapt & \CsThreeUnkPPT & \CsThreeUnkColl & \CsThreeUnkGap \\
\bottomrule
\end{tabular}}
\end{table}

\paragraph{Results.} Table~\ref{tab:dcp} has three findings.

\emph{Known secret, recorded pure instance: no collective advantage over the best local strategy.} A two-step adaptive circuit measures the first wire in the $XY$ plane and then the second wire in a direction chosen by the first outcome. With this one bit of feed-forward it reaches the collective optimum $\CsModKnownColl$ to floating-point precision, so here the PPT bound coincides with the collective value and with an achieved local score. This agrees with the known theorem that local operations and classical communication (LOCC) reach the Helstrom bound for any two pure multipartite states~\cite{VirmaniSacchiPlenioMarkham2001}; both hypotheses here are pure product states. The collective measurement does beat the two product rules the note evaluated, the $X$-basis product ($\CsModKnownX$) and the product of local Helstrom readouts ($\CsModKnownLocHel$), and the best product receiver the search found with free local directions reaches $\CsModKnownProd$. Against the best local strategy, the instance shows no collective advantage.

\emph{Mixed packets: a small certified local gap on the three-packet instances.} With a known secret, the three-packet collective optimum is $\CsThreeKnownColl$, while no local protocol of any kind exceeds the PPT bound $\CsThreeKnownPPT$. That is a certified gap of at least $\CsThreeKnownGap$; the adaptive search reaches $\CsThreeKnownAdapt$, leaving an open window of at most $\CsThreeKnownWindow$ for richer local protocols. With unknown sign the certified gap is at least $\CsThreeUnkGap$, and the adaptive search is within $\CsThreeUnkWindow$ of the PPT bound. The two-packet mixed instances behave differently. With a known secret the adaptive search reaches the collective optimum $\CsTwoKnownColl$ and the certified gap is $\CsTwoKnownGap$. With unknown sign the certified gap is at least $\CsTwoUnkGap$ but below $\CsTwoUnkGapUp$, and the open window is at most $\CsTwoUnkWindow$. A certified gap near the three-packet size is therefore not a property of every mixed instance.

\emph{Unknown sign, recorded instance: open.} Sign symmetrization makes the alternative hypothesis a mixture. The product and adaptive searches find nothing above the $X$ basis ($\CsModUnkAdapt$), and the PPT bound does not fall below the collective score $\CsModUnkColl$. The PPT program's primal solution is, up to the solver's feasibility tolerance, a PPT measurement reaching $\CsModUnkPrimal$, against a collective score of $\CsModUnkCollLong$. On two qubits a PPT effect is separable~\cite{Horodecki1996}, so a separable measurement comes within solver precision of the collective score; this rests on the solver's primal solution, not on a certificate. Whether an LOCC protocol does as well is undecided. The best local protocol lies between $\CsModUnkAdapt$ and $\CsModUnkPPT$, an open window of at most $\CsModUnkWindow$.

\paragraph{Scope.} The bridge note says: ``The previously stated universal strict three-tier hierarchy is FALSE.'' and ``The two specified product strategies have no universal ordering.'' and ``This also does not compare against the optimal adaptive/separable receiver or solve unknown-secret DCP.'' The circuit-search report calls itself ``Not a registry entry; for the research agents to evaluate.'' Its documentation adds: ``Nothing here is asymptotic.'' and ``Only the PPT column and the collective column are bounds.'' and, for the unknown-sign recorded instance, ``Whether LOCC reaches it is undecided: the search found nothing above the X basis.'' All instances have at most $\DcpThreePackets$ packets, and a binary test with a known or sign-ambiguous secret is not a DCP solver.
\begin{flushleft}
The bridge report sets \texttt{accepted\_speedup\_candidate}, \texttt{novelty\_claim} and \texttt{universal\_strict\_three\_tier\_hierarchy\_claimed} to \texttt{false}, and \texttt{proof\_status} to \texttt{LOCAL\_DERIVATION\_REVIEW\_PENDING\_NOT\_FORMALLY\_VERIFIED}.
\end{flushleft}

%--------------------------------------------------------------------
\section{Result C: a classical sampler for the polar Gibbs source}
\label{sec:polar}

\paragraph{Setting.} Let $E$ have $N$ independent Bernoulli$(p)$ bits with $p=(1-c)/2$, where $N$ is a power of two and $c$ a rational overlap. The polar transform gives $U=G(E)$. A public frozen set $F$ of polar coordinates defines parity checks $H_F$, so that $U_F = H_F E$. Given public clause signs $J$, let $s = H_F J$, draw $E$ conditioned on $U_F = s$, and set $y = E \oplus J$. Then $H_F y = 0$ and
\[
\operatorname{Pr}(y) \;=\; \frac{1}{Z_J} \prod_{j} \left(1 + c\,(-1)^{y_j + J_j}\right),
\]
with $Z_J$ the normalizing constant over the kernel. This is the signed parity-Ising Gibbs law on the kernel of $H_F$. Because $H_F$ has full row rank, uniformly random signs $J$ give a uniformly random syndrome $s$. All matrices and probabilities are explicit; no oracle is supplied. The instances use $c=\PolarOverlap$, so $p=\PolarSourceP$. The \emph{rate} of a design is $|F|/N$, the rate of the code that the DQI reduction must decode; the kernel code itself has rate $1-|F|/N$.

\paragraph{Claim tested.} Zhou, Shutty, Sellke and Jordan relate quantum decoding to Gibbs sampling when the clause signs are uniformly random, and bound the expected total variation~\cite{ZhouShuttySellkeJordan2026}. Polar successive cancellation with quantum messages gives a natural coherent decoder whose message graphs are trees~\cite{PiveteauRenes2021,MandalBrandsenPfister2024}; related convergence results for random LDPC codes are in~\cite{MandalPiveteauRenesPfister2026}. Piveteau, Quek and Franco Garrido replace the coherent classical decoder in this reduction with quantum decoders over the same pure-state channels. For commuting Hamiltonians they dequantize their unambiguous-discrimination decoder into a classical rejection sampler with the same temperature threshold, and they report a BPQM threshold beyond the classical Shannon limit on dilute spin glasses~\cite{PiveteauQuekFrancoGarrido2026}. The proposal was to use such a polar receiver inside the DQI Gibbs reduction to sample this law faster than classical methods. The project's earlier pass left one blocker: its error criterion did not construct a suitable public frozen set without an exponential source table.

\paragraph{Classical competitor.} SC source shaping walks $U_0,\ldots,U_{N-1}$, forces frozen coordinates to their syndrome values, and samples the others from the exact unrestricted conditional $P(U_i \mid U_{<i})$, using exact rational arithmetic and $O(N^2)$ rational operations in the reference implementation. It ignores future frozen constraints, so in general it is not exact for a specific syndrome; it is exact for a contiguous frozen prefix. By a result of Korada and Urbanke~\cite{KoradaUrbanke2010}, the joint law with uniform frozen values is within total variation $\tfrac12\sum_{i\in F} m_i$ of the source law. Passing to the conditional law given the syndrome at most doubles this on average over uniform signs:
\[
\mathbb{E}_J\, \TV \;\le\; \sum_{i \in F} m_i, \qquad m_i = \mathbb{E}\,\lvert P(U_i = 0 \mid U_{<i}) - P(U_i = 1 \mid U_{<i}) \rvert .
\]
Each $m_i$ equals the unsquared root fidelity of the reversed synthesized classical-quantum bit channel, an instance of channel/code duality~\cite{Renes2017,RengaswamyPfister2021}; this identity is a local derivation, checked numerically only at the smallest lengths and marked review pending. Hence any polar Gibbs proposal justified by a vanishing sum of these root fidelities also supplies this classical average-sign sampler. The average is essential: on a length-$\CalN$ calibration with $p=\CalNoise$, the uniform-syndrome average error is $\CalAvg$ but the worst syndrome has error $\CalWorst$.

To make the competitor constructive, the project certifies upper bounds on every $m_i$ without a source table. Overlaps are kept on the nonuniform mesh $(j/G)^{r}$, $j=0,\ldots,G$, with dyadic herald masses of $\PolarMassBits$-bit precision. Each channel combination is computed exactly in integers, overlaps are rounded up and masses down, and discarded mass is sent to an uninformative flag. The source note argues that the posterior-magnitude recursion is a combination of flagged pure-state channels and that each of these roundings is a physical degradation of such a channel, under which root fidelity cannot decrease; the computed means are then upper bounds on the $m_i$. Both steps are local derivations marked review pending, so every certified bound below is conditional on them. The frozen set is chosen as the coordinates with the smallest certified bounds. Bounded-alphabet polar construction has substantial prior art~\cite{TalVardy2011}, and SC source shaping builds on source polarization~\cite{Arikan2010Source,HondaYamamoto2013}.

\begin{table}[t]
\centering
\small
\caption{Result C, certified designs at $c=\PolarOverlap$, $p=\PolarSourceP$. ``Largest set'' is the largest frozen set whose certified average-sign error stays within $\PolarTarget$, for this bound vector. ``Half-rate bound'' is the certified bound for a frozen set of half the length. Sources: \texttt{polar\_source\_design.json} and \texttt{polar\_source\_design\_n32768\_g256\_p3.json}.}
\label{tab:polar}
\begin{tabular}{rllll}
\toprule
$N$ & Mesh $(G, r)$ & Largest set within $\PolarTarget$ & Its rate & Half-rate bound \\
\midrule
\CoarseNA & (\CoarseGrid, \CoarsePower) & \CoarseFA & \CoarseRateA & \CoarseHalfA{} (trivial) \\
\CoarseNB & (\CoarseGrid, \CoarsePower) & \CoarseFB & \CoarseRateB & \CoarseHalfB{} (trivial) \\
\CoarseNC & (\CoarseGrid, \CoarsePower) & \CoarseFC & \CoarseRateC & \CoarseHalfC{} (trivial) \\
\CoarseND & (\CoarseGrid, \CoarsePower) & \CoarseFD & \CoarseRateD & \CoarseHalfD{} (trivial) \\
\PolarN & (\PolarGrid, \PolarPower) & \PolarMaxF & \PolarMaxRate & Eq.~(\ref{eq:polartv}) \\
\bottomrule
\end{tabular}
\end{table}

\paragraph{Results.} Table~\ref{tab:polar} lists the designs. On the coarse mesh the half-rate targets fail certification (their bound is the trivial value $\CoarseHalfD$), and the report labels them unknown at that resolution, not hard. On the finer mesh at $N=\PolarN$, the half-rate design with $|F|=\PolarF$ is certified:
\begin{equation}
\label{eq:polartv}
\mathbb{E}_J\, \TV \;\le\; \PolarTV \;\le\; \PolarTVDec .
\end{equation}
The first bound is the exact rational stored in the report; the second rounds it up. Like every certificate in this section, it is conditional on review-pending local derivations: the Gibbs mapping and hybrid bound above, and the two steps behind the upper bounds on the $m_i$ (the pure-state channel correspondence and the degradation argument). The construction used \PolarCombos{} channel combinations and \PolarVisits{} active support-pair visits, and this setup cost is disclosed rather than treated as free advice. At the same length the largest frozen set within $\PolarTarget$ has $\PolarMaxF$ constraints. A separate rate-$\EighthRate$ design at $N=\EighthN$ is also certified, but its rate was chosen after exploration and was not declared in advance. The reports also note that both certified rates, $\PolarRate$ and $\EighthRate$, exceed $1-c$, the rate ceiling of a different mechanism the project analysed earlier (product unambiguous state discrimination followed by Prange-style decoding); that ceiling applies to that mechanism only, not to all classical algorithms. A classical sampler matched to unambiguous discrimination also appears in~\cite{PiveteauQuekFrancoGarrido2026}. For this polar family and the uniform-sign average, then, the matching classical route is constructive.

\paragraph{Scope.} Polar source design note: ``No novel quantum algorithm, all-code no-go or matched runtime advantage.'' and ``It does NOT refute their random-hypergraph results: our code family is polar.'' and ``The certificate is not a quantum advantage.'' and ``It does NOT prove a fixed-grid asymptotic convergence theorem, a general Gibbs sampler, or an all-sign result.'' and ``An adversarial sign regime requires a different proof, not this average bound.'' The same note records that a Gibbs draw at the largest length has not been executed. Polar conditional sampler note: ``There is NO claim that every adversarial sign pattern is handled well.'' and ``That is a CONDITIONAL, review-pending application of known mechanisms, not a worst-case-sign no-go, not all codes, and not a result established by the small numerics.'' BPQM Gibbs frontier note: ``No novel quantum algorithm, finite-block receiver, classical-hardness theorem, or Shor-level speedup is claimed.''
\begin{flushleft}
The large report sets \texttt{accepted\_speedup\_candidate}, \texttt{all\_codes\_dequantized}, \texttt{worst\_case\_sign\_sampler\_certified} and \texttt{quantum\_decoder\_executed} to \texttt{false}, and \texttt{guarantee\_is\_average\_uniform\_sign\_not\_worst\_sign} and \texttt{local\_degradation\_and\_hybrid\_derivation\_review\_pending} to \texttt{true}.
\end{flushleft}

%--------------------------------------------------------------------
\section{Rediscovery audit of the negative-results registry}
\label{sec:lit}

The project audited its own negative results for rediscovery. The $\RedTotal$ registry records were clustered by shared reasoning. For each cluster a judge read some of the members, stated the shared barrier in standard terms, searched arXiv and assigned a label: \emph{published} (a retrieved abstract states the same barrier for the same setting), \emph{adjacent} (known technique or barrier in a different setting, or only implied), \emph{not-found}, or \emph{unclear}. A label needs evidence: a \emph{published} label ``requires a citation whose abstract contains the judge's quote verbatim.'' The judge could override the cluster label for a member it read. An unread member inherits the cluster label only when at least half of the read members kept that label and one of those shares the member's track; otherwise it is withheld and counted as unclear. $\RedShown$ records were read by a judge, $\RedPropagated$ inherited their cluster's label and $\RedWithheld$ were withheld.

\begin{table}[t]
\centering
\small
\caption{Rediscovery audit: records per label, overall, among the records a judge read, and by track (all records). Withheld records count as unclear. A dash means the report has no record with that label in that row. Source: \texttt{research/audits/rediscovery\_audit.json}.}
\label{tab:audit}
\begin{tabular}{lrrrr}
\toprule
Track & Published & Adjacent & Not found & Unclear \\
\midrule
All records (of \RedTotal) & \RedPublished & \RedAdjacent & \RedNotFound & \RedUnclear \\
Read by a judge (of \RedShown) & \RedReadPublished & \RedReadAdjacent & \RedReadNotFound & \RedReadUnclear \\
\midrule
Code equivalence & \RedCodePub & \RedCodeAdj & \RedCodeNf & \RedCodeUnc \\
Coset observables / nonabelian HSP & \RedCosetPub & \RedCosetAdj & \RedCosetNf & \RedCosetUnc \\
Dihedral coset / hidden shift & \RedDcpPub & \RedDcpAdj & \RedDcpNf & \RedDcpUnc \\
Legacy toy oracles & \RedToyPub & \RedToyAdj & \RedToyNf & \RedToyUnc \\
Query lower bounds & \RedQueryPub & \RedQueryAdj & \RedQueryNf & \RedQueryUnc \\
Other & \RedOtherPub & \RedOtherAdj & \RedOtherNf & \RedOtherUnc \\
Applied-QC barriers from the literature & \RedAppliedPub & \RedAppliedAdj & \RedAppliedNf & \RedAppliedUnc \\
Lattices & \RedLatticePub & \RedLatticeAdj & \RedLatticeNf & \RedLatticeUnc \\
DQI / OPI / prime characters & \RedDqiPub & \RedDqiAdj & \RedDqiNf & \RedDqiUnc \\
\bottomrule
\end{tabular}
\end{table}

\paragraph{Headline and per-track shares.} Table~\ref{tab:audit} gives the counts; the report stores counts, and we quote shares as counts. $\RedPublished$ of $\RedTotal$ records restate a published barrier, $\RedAdjacent$ are adjacent to known work, $\RedUnclear$ are unclear (including the $\RedWithheld$ withheld records), and only $\RedNotFound$ came back not-found. Only the labels on records a judge read rest on a reading of the record itself: of those $\RedShown$ records, $\RedReadPublished$ were labelled published, $\RedReadAdjacent$ adjacent, $\RedReadNotFound$ not-found and $\RedReadUnclear$ unclear; the remaining labels were inherited ($\RedPropagated$) or withheld ($\RedWithheld$). Among the three largest tracks, the published share is largest for coset observables / nonabelian HSP ($\RedCosetPub$ published, $\RedCosetAdj$ adjacent, $\RedCosetUnc$ unclear), next for code equivalence ($\RedCodePub$, $\RedCodeAdj$, $\RedCodeUnc$) and smallest for the dihedral coset track ($\RedDcpPub$, $\RedDcpAdj$, $\RedDcpUnc$). Among the smaller tracks, the applied-QC barriers taken from the literature have $\RedAppliedPub$ published and $\RedAppliedUnc$ unclear records; the legacy toy-oracle, lattice and DQI/OPI tracks have none published, and the query-lower-bound track has $\RedQueryPub$.

\paragraph{Agreement checks.} Before judging, $\RedGroundN$ records restating textbook barriers were fixed as ground truth; $\RedGroundPub$ of them came back published. This is a weak check: the judge instructions used a barren-plateau query and citation as their worked examples, author--year references in record text (such as McClean et al.\ 2018 in the barren-plateau record) stayed visible because only arXiv identifiers were hidden, and most of the published ground-truth hits are single-record overrides inside one mixed cluster that is labelled unclear as a whole. A fresh reviewer re-read $\RedSpotSampled$ randomly chosen judged records with their labels and citations. The labels of $\RedSpotStale$ of them changed after that review, when the propagation guard withheld them; the reviewer had agreed with $\RedSpotStaleAgree$ of those labels as they were then. On the $\RedSpotN$ whose label is unchanged the reviewer agreed with $\RedSpotAgree$: $\RedSpotShownAgree$ of $\RedSpotShownN$ judge-read records and $\RedSpotPropAgree$ of $\RedSpotPropN$ inherited ones.

\paragraph{The audit's own caveats.} ``Judging is done by in-session subagents.'' ``Labels are cluster-level. The reviewer found records whose own barrier differs from their cluster's, including one that was mis-clustered.'' ``Unclear is mostly templated records (``Falsified or refined by exact theorem evaluation'') and mixed clusters.'' The ground-truth records that were not labelled published (lossy boson sampling, Gaussian continuous-variable circuits, QAOA locality) sit in a mixed cluster labelled unclear, and one ground-truth hit carries its cluster's citation rather than the papers specific to it.

\paragraph{Relation to Results A--C.} None of the three results is a registry record, and the audit's DQI/OPI track holds a single record labelled unclear, so the audit does not measure their novelty. The source notes name the prior art themselves. In Result A, the norm and factorization machinery, the arithmetic-progression statistics and the trace-function bounds are standard~\cite{KuoLiuTotani2026,BankBarySorokerRosenzweig2013,FouvryKowalskiMichel2014}, and the quantum side is the known DQI mechanism~\cite{JordanDQI2024,SunWootters2026}, whose margin over Prange completion under Hamming decoding is bounded in~\cite{MazumderNiroula2026}. Diagonal-equation algorithms~\cite{IvanyosSantha2015} are noted in the OPI note as covering a different regime. For the known-secret pure instance of Result B the decisive local result is a published theorem~\cite{VirmaniSacchiPlenioMarkham2001}; the mixed and unknown-sign findings rest on the PPT program. In Result C, polarization, SC source shaping, the hybrid bound and channel duality are known~\cite{Arikan2009,Arikan2010Source,HondaYamamoto2013,KoradaUrbanke2010,TalVardy2011,Renes2017,RengaswamyPfister2021}. What the project adds is finite, fixed-instance evaluation with explicit non-claims, not new algorithms.

%--------------------------------------------------------------------
\section{Limitations}
\label{sec:limits}

\begin{itemize}
\item \textbf{Review status.} The source notes mark the local derivations (Morse monodromy, exterior trace bound, finite Gram certificates, polar degradation and hybrid bounds, the DCP bridge) as review pending. None has had independent or formal review.
\item \textbf{Finite instances.} All three results are finite. Result A uses saved seeded inputs, and seeded replays are not independent-input confidence statements. Result B has at most $\DcpThreePackets$ packets. Result C certifies specific lengths and meshes and proves no asymptotic convergence.
\item \textbf{Result A.} The classical high-rate witnesses are for the saved inputs. For the saved $N=\OpiN$ inputs, the fixed-input lower bounds in the original field are zero, which shows only that these conservative bounds do not apply there; the positive bound was obtained for one input re-posed over larger declared primes, a different problem. At the low rate there is no classical match for the conditional quantum certificate and no implemented quantum decoder, so no conclusion is drawn either way. Not every saved high-rate input has a perfect classical witness. One quadratic input remains imperfect: the mixed input with $N=\OpiMidN$, $K=\OpiMidK$, $r=\OpiMidR$ and seed $\OpiMidSeed$, over the same $\OpiMidPrimeBits$-bit prime, reaches $\OpiMidMixed$ of $\OpiMidN$ under power completion and adaptive repair, and the norm search was not run on the $N=\OpiMidN$ inputs. Saved high-rate mixed inputs with higher power characters also remain imperfect: $\CubeMixedA$ and $\CubeMixedB$ of $\CubeN$ for $r=\CubeR$ ($K=\CubeK$), and $\QuartMixedA$ and $\QuartMixedB$ of $\QuartN$ for $r=\QuartR$ ($K=\QuartK$). Only power completion and a linear baseline were run on these, and the $r=\QuartR$ inputs are over $p=\QuartPrime$. Other orientation patterns and other fields are not settled.
\item \textbf{Result B.} ``Best product'' and ``best adaptive'' are search results, so they are lower bounds on their optima. Richer local protocols (multi-outcome POVMs per wire, two-way rounds) may close the open windows. The PPT gaps are small in total-variation distance, and a binary test is not a DCP algorithm.
\item \textbf{Result C.} The guarantee is an average over uniform signs. Adversarial signs, other code families (random LDPC, cyclic) and non-SC joint measurements are not covered. The setup cost at the largest length is substantial, and no quantum receiver was compiled for comparison.
\item \textbf{No matched resource comparisons.} None of the three results compares a quantum and a classical procedure at matched runtime or gate cost.
\item \textbf{Audit.} The audit labels come from model judges, mostly at cluster level; the spot check is small, the ground truth is a weak check (Section~\ref{sec:lit}), the failure-mode agreement measures one model's repeatability, and the audit does not cover Results A--C.
\end{itemize}

%--------------------------------------------------------------------
\section{Reproducibility}
\label{sec:repro}

Every number in this paper is a macro defined in \url{numbers.tex}. That file is generated by \url{make_numbers.py} from \url{numbers_map.json}, which names, for each macro, a committed JSON report, a path into it and a format. A path that no longer resolves stops \url{make_numbers.py}; the only exception is an audit count declared optional, which prints as a dash when the report omits it, as the report does for zero counts. A value that changes does not stop anything by itself: \url{numbers.tex} stays stale until the script is rerun, and \url{tests/tools/test_paper_numbers.py} fails while the committed \url{numbers.tex} differs from what the script would generate. Running \texttt{make_numbers.py --check} also lists any command in this file that is neither a generated number nor a known LaTeX command.

The script \texttt{reproduce.sh} in the paper directory and the \texttt{qsearch-classical} package (\url{packages/qsearch_classical}) regenerate the numbers. The package vendors the prime-character modules (OPI, norm attack, Morse attack, exterior bound, capacity source, finite readout) and the polar modules (conditional sampler, source design), each writing its report under \url{research/classical_baselines/}. Its tests recompute the OPI, norm-attack, capacity-source, finite-readout, conditional-sampler and small source-design reports exactly, ignoring source hashes and library versions; the package notes that the Morse and exterior-bound runs are too slow for continuous integration and were checked once by hand. \texttt{reproduce.sh} runs these modules from the vendored source, not from an installed copy, and stops at once if one of the packages it needs (python-flint, sympy, numpy, scipy, cvxpy) is missing. It reruns all of these modules into a fresh directory, including the large polar design, reruns the DCP circuit search (\url{tools/circuit_search/dcp_readout.py}), and then runs \texttt{make_numbers.py --compare} to compare every macro with its committed value. Three sources are not regenerated by it: the bridge-note values of Result B (\url{theorems/dcp_multi_packet_collective_helstrom_product_likelihood_bridge.py}) and the two audits, which \url{tools/audit/rediscovery.py} and \url{tools/audit/failure_modes.py} rebuild from their stored judgements without re-judging. A full run of the script on a laptop reproduced every regenerated number exactly; the remaining macros come from the three sources above.

\bibliographystyle{plain}
\bibliography{references}

\end{document}
